\documentclass[11pt]{article}
\usepackage[margin=1in]{geometry}
\usepackage{amsmath,amssymb,amsthm,mathtools}
\usepackage{booktabs}
\usepackage{enumitem}
\usepackage{hyperref}
\hypersetup{colorlinks=true,linkcolor=blue,citecolor=blue,urlcolor=blue}

\theoremstyle{plain}
\newtheorem{proposition}{Proposition}
\newtheorem{lemma}{Lemma}
\theoremstyle{definition}
\newtheorem{definition}{Definition}
\newtheorem{construction}{Construction}
\newtheorem{conjecture}{Conjecture}
\theoremstyle{remark}
\newtheorem{result}{Computational Result}
\newtheorem{remark}{Remark}

% --- certainty-tier tags, used throughout ---
\newcommand{\Proven}{\textsf{[proven]}}
\newcommand{\Symbolic}{\textsf{[symbolic, all $m_D$]}}
\newcommand{\Measured}{\textsf{[computer-verified]}}
\newcommand{\Conj}{\textsf{[conjecture]}}
\newcommand{\Interp}{\textsf{[interpretation]}}

\newcommand{\Sym}{\operatorname{Sym}}
\newcommand{\Gal}{\operatorname{Gal}}
\newcommand{\Ftil}{\widetilde{F}}
\newcommand{\C}{\mathbb{C}}
\newcommand{\Q}{\mathbb{Q}}
\newcommand{\Z}{\mathbb{Z}}

\title{\textbf{Imprimitive Wreath Monodromy of Self-Composed\\ Constraint Surfaces}\\[4pt]
\large The seam as imprimitivity system, the omega substrate,\\ and a forward/backward recovery asymmetry}
\author{William Lloyd\\ \small Independent researcher, Coffs Harbour NSW, Australia\\[2pt]
\small \emph{Draft --- session results; certainty tiers marked inline.}}
\date{\today}

\begin{document}
\maketitle

\begin{abstract}
We study the branched-cover structure that arises when a constraint surface
$F(D,S,P)=0$, carrying a directive--substrate--product (DBP) role assignment, is
composed with itself by identifying the product of one stage with the directive
of the next. Eliminating the shared variable (the \emph{seam}) yields a
four-variable hypersurface $\Ftil(D_1,S_1,S_2,P_2)=0$ whose monodromy we compute
directly. Three findings stand out. First, there is no single monodromy group:
each role-output cover carries its own, and each is a \emph{full wreath product
that is imprimitive with respect to the seam}. For $F = D^{m_D}+DS+P^{m_P}-3$ the
product cover is $\Sym(m_P)\wr\Sym(m_P)$ and, at $m_P=2$, the directive cover is
$\Sym(m_D)\wr\Sym(m_D)$, with degree spectrum $(m_D^2,m_D,m_P,m_P^2)$. Second, the
second-stage substrate $S_2$ is an \emph{omega role}: it controls the position of
the downstream product branch while leaving the seam invariant
($\partial A/\partial S_2=0$, $\partial C/\partial S_2=1$), and looping it steers
exactly one product branch pair at a time. Third, iterated composition produces a
sharp recovery asymmetry: reconstructing the root directive over $n$ stages is
degree $m_D^{\,n}$ while producing the final product is degree $m_P^{\,n}$. We
also bound a tempting ``wave/phase'' reading: the abelianisation of the directive
group is $(\Z/2)^2$, so its $U(1)$ characters are signs and the higher (cubic)
structure is non-abelian and phaseless. Results are computer-verified on two
surfaces, the cubic ($m_D=3,m_P=2$) and the keystone ($m_D=m_P=2$); the keystone
was chosen as an adversarial falsification test and passed.
\end{abstract}

\tableofcontents

\section{Introduction}

A directive-based product (DBP) presents an algebraic relation
$D\oplus S=P$ as a constraint surface $F(D,S,P)=0$ together with an assignment of
\emph{roles}: $D$ is the directive (the active operator), $S$ the substrate (the
passive operand), and $P$ the product. The framework's defining intuition is an
asymmetry: $D$ acting on $S$ yields $P$, whereas ``$S$ acting on $D$'' is a
category error. As an \emph{algebraic} object, however, $F=0$ is symmetric---one
may solve for any variable---so the directional restriction appears to dissolve.
The first contribution of this paper is to show that it does not dissolve; it is
\emph{re-encoded in the geometry}, as an asymmetry in the degree and monodromy of
the inverse problems.

The natural way to probe this is composition. We compose a DBP surface with itself
by the \emph{forward glue} $D_2 := P_1$: the product of stage one becomes the
directive of stage two (output feeds next operator). Eliminating the shared
variable produces a single hypersurface $\Ftil=0$, and we ask what survives of the
role structure in its branched-cover geometry.

\paragraph{Summary of results and certainty tiers.}
Throughout we tag each statement by how strongly it is established: \Proven\ (an
exact algebraic fact), \Symbolic\ (verified symbolically for the whole family
$m_D=2,3,4,\dots$ at $m_P=2$ but not proved in closed form for all $m_D$),
\Measured\ (established by numerically certified monodromy on a specific surface,
with the gates of \S\ref{sec:methods}), \Conj\ (a conjecture, with the supporting
evidence stated), and \Interp\ (an interpretive reading, not a theorem).

\begin{enumerate}[leftmargin=1.4em]
\item \textbf{No single Galois group; per-cover wreaths.} Each role-output cover of
$\Ftil$ has its own monodromy group, and each is a full wreath product imprimitive
with respect to the seam (\S\ref{sec:product}--\ref{sec:directive}). \Measured/\Symbolic
\item \textbf{Degree spectrum law.} For $F=D^{m_D}+DS+P^{m_P}-3$ at $m_P=2$, the
spectrum of $(D_1,S_1,S_2,P_2)$ in $\Ftil$ is $(m_D^2,m_D,m_P,m_P^2)$. \Symbolic
\item \textbf{Product cover $\cong \Sym(m_P)\wr\Sym(m_P)$}, directive cover
$\cong\Sym(m_D)\wr\Sym(m_D)$ at $m_P=2$; the seam is the block system in both.
Cubic: $D_4$ (product) and $S_3\wr S_3$ of order $1296$ (directive). Keystone:
$D_4$ on both. \Measured
\item \textbf{$D_4$ equals the seam-partition stabiliser.} For the product cover,
the monodromy group is exactly the stabiliser in $S_4$ of the seam-aligned
partition $\{\{D_1,S_1\},\{S_2,P_2\}\}$; three independent derivations agree
(\S\ref{sec:product}). \Measured
\item \textbf{The omega substrate.} $S_2$ controls the downstream product-branch
position with the signature $\partial A/\partial S_2=0$, $\partial C/\partial
S_2=1$; loops in $S_2$ steer exactly one product pair (\S\ref{sec:omega}). \Measured
\item \textbf{Recovery asymmetry.} Over $n$ stages, root-directive recovery is
degree $m_D^{\,n}$ and final-product production is degree $m_P^{\,n}$; for the cubic
$3^n$ vs $2^n$ (\S\ref{sec:iterate}). \Symbolic
\item \textbf{Abelian boundary on the phase reading.} The abelianisation of
$S_3\wr S_3$ is $(\Z/2)^2$; its $U(1)$ characters take values in $\{\pm1\}$, and the
cube-root structure lives in the non-abelian part with no $U(1)$ character
(\S\ref{sec:abelian}). \Measured
\end{enumerate}

\section{The self-glue construction}
\label{sec:construction}

\begin{definition}[DBP surface, worked family]
We work with
\[
  F(D,S,P) \;=\; D^{m_D} + D\,S + P^{m_P} - 3 \;=\; 0,
\]
with integers $m_D,m_P\ge 2$, role assignment $(D,S,P)=(\text{directive},
\text{substrate},\text{product})$. The two running examples are the
\emph{cubic} ($m_D=3,m_P=2$, i.e.\ $D^3+DS+P^2-3$) and the \emph{keystone}
($m_D=m_P=2$, i.e.\ $D^2+DS+P^2-3$).
\end{definition}

\begin{construction}[Forward self-glue]
\label{con:glue}
Take two copies of $F$. Stage one is $F(D_1,S_1,P_1)=0$; stage two is
$F(D_2,S_2,P_2)=0$ with the \emph{seam identification} $D_2:=P_1$. Writing
\[
  A \;:=\; 3 - D_1^{m_D} - D_1 S_1, \qquad\text{so that}\qquad P_1^{m_P}=A,
\]
stage two becomes $P_1^{m_D}+P_1 S_2 + P_2^{m_P}-3=0$. Eliminating the seam
variable $P_1$ (by resultant with $P_1^{m_P}-A$) yields a four-variable
hypersurface
\[
  \Ftil(D_1,S_1,S_2,P_2) = 0 .
\]
\end{construction}

\noindent At $m_P=2$ the seam relation is $P_1^2=A$ and the elimination is
explicit. For the cubic, with $C:=A+S_2$ and $B:=3-P_2^2$, stage two reads
$P_1 C = B$, so $P_1^2 = A$ gives
\[
  \boxed{\;\Ftil_{\text{cubic}} = B^2 - A\,C^2\;}\qquad (A = 3-D_1^3-D_1S_1,\; C=A+S_2,\; B=3-P_2^2).
\]
For the keystone, $P_1^2=A$ makes the directive term $D_2^2=P_1^2=A$ a constant in
$P_1$, so stage two reads $P_1 S_2 = 3-A-P_2^2$ and hence
\[
  \boxed{\;\Ftil_{\text{key}} = (3-A-P_2^2)^2 - A\,S_2^2\;}\qquad (A=3-D_1^2-D_1S_1).
\]

\begin{definition}[Lifted correspondence]
Rather than work on $\Ftil$ alone we retain the seam $P_1$ as an explicit fibre
coordinate. At $m_P=2$ the lifted system is the tower of two quadratics
\[
  P_1^2 = A, \qquad P_2^2 = 3 - P_1\,(\text{stage-two coupling}),
\]
which for the cubic is $P_2^2 = 3 - P_1 C$ and for the keystone $P_2^2 = 3 -
P_1(P_1+S_2)$. The seam is the intermediate sheet over which the product sits.
\end{definition}

\section{Role-output covers and the degree spectrum}
\label{sec:spectrum}

$\Ftil=0$ is a three-dimensional hypersurface in four variables; fixing which
three variables form the base and solving for the fourth gives a finite branched
cover. We call these the \emph{role-output covers}.

\begin{proposition}[No single Galois group]
\label{prop:nosingle}
The monodromy group of $\Ftil$ depends on the choice of output role. Concretely,
the directive cover (solve for $D_1$) and the product cover (solve for $P_2$) have
different degrees and different groups. A naive ``$\Gal(\Ftil/\text{base})$'' over
a base in which the seam variable $A$ is treated as a free parameter is the group
of a \emph{decoupled} problem and does not describe any actual cover.
\end{proposition}

\begin{remark}
Proposition~\ref{prop:nosingle} is a correction of a natural error: one is tempted
to assign $\Ftil$ a single group (e.g.\ a product $S_3\times D_4$ obtained by
treating $A$ as free). This fails on a counting obstruction---the directive cover
is transitive of degree $9$, so $9$ must divide its order, but
$|S_3\times D_4|=48$ is not divisible by $9$. The honest object is the
per-output-cover group.
\end{remark}

\begin{result}[Degree spectrum law, $m_P=2$]
\label{res:spectrum}
For $F=D^{m_D}+DS+P^2-3$ the degrees of $(D_1,S_1,S_2,P_2)$ in $\Ftil$ are
\[
  (\deg_{D_1},\deg_{S_1},\deg_{S_2},\deg_{P_2}) \;=\; (m_D^2,\; m_D,\; m_P,\; m_P^2),
\]
verified symbolically for $m_D\in\{2,3,4\}$. \Symbolic\quad The directive and its
own substrate carry $m_D$-powers; the product and its substrate (the omega role,
\S\ref{sec:omega}) carry $m_P$-powers.
\end{result}

\begin{center}
\begin{tabular}{lcccc}
\toprule
surface & $\deg_{D_1}$ & $\deg_{S_1}$ & $\deg_{S_2}$ & $\deg_{P_2}$ \\
\midrule
cubic $(m_D{=}3,m_P{=}2)$ & $9$ & $3$ & $2$ & $4$ \\
keystone $(m_D{=}2,m_P{=}2)$ & $4$ & $2$ & $2$ & $4$ \\
generic $m_D$ (at $m_P{=}2$) & $m_D^2$ & $m_D$ & $2$ & $4$ \\
\bottomrule
\end{tabular}
\end{center}

\noindent The keystone spectrum $(4,2,2,4)$ is symmetric under $D_1\leftrightarrow
P_2$ precisely because $m_D=m_P$; the cubic cannot be, because $m_D\neq m_P$. This
$D_1\leftrightarrow P_2$ symmetry is the algebraic shadow of the directive/product
duality.

\section{The product cover}
\label{sec:product}

Solving $\Ftil$ for $P_2$ over $(D_1,S_1,S_2)$: $A$ is single-valued, the seam
$P_1$ has $m_P$ values ($P_1^{m_P}=A$), and the product $P_2$ has $m_P$ values.
Both layers are $m_P$-fold extractions, so the cover factors as $m_P$ seam blocks
of $m_P$ sheets each.

\begin{result}[Product-cover monodromy]
\label{res:product}
The product cover is the full wreath $\Sym(m_P)\wr\Sym(m_P)$, of degree $m_P^2$,
imprimitive with the seam ($\pm P_1$ for $m_P=2$) as block system. \Measured\ at
$m_P=2$. For both the cubic and the keystone this group is
\[
  \Sym(2)\wr\Sym(2) \;\cong\; D_4 \quad(\text{order } 8).
\]
\end{result}

\noindent The measurement uses three based loops from the basepoint $A_0=1$,
$C_0=1$ (cubic) on the four lifted sheets
$0{=}({+}P_1,{+}P_2)$, $1{=}({+}P_1,{-}P_2)$, $2{=}({-}P_1,{+}P_2)$,
$3{=}({-}P_1,{-}P_2)$. With residual gates on \emph{both} equations
($\le 4\times10^{-15}$), bijective closure, and motion-ratio safety:
\[
  \tau_A=(0\,2)(1\,3)\ \text{(seam)},\quad
  \tau_+=(0\,1)\ \text{($+P_1$ product)},\quad
  \tau_-=(2\,3)\ \text{($-P_1$ product)}.
\]
The conjugation $\tau_A\tau_+\tau_A^{-1}=\tau_-$ holds, so the group is
non-abelian; it contains the $4$-cycle $(0\,2\,1\,3)$ and has order $8$.

\begin{result}[$D_4$ is the seam-partition stabiliser]
\label{res:stab}
The product-cover group equals the stabiliser in $S_4$ of the seam-aligned
partition $\{\{D_1,S_1\},\{S_2,P_2\}\}$ (equivalently $\{\{0,1\},\{2,3\}\}$ on the
lifted sheets), the order-$8$ ``$V_4$-core''. \Measured
\end{result}

\noindent Three independent derivations agree on this order-$8$ group: (i) the
degree spectrum $(9,3,2,4)$ is faithful with four distinct degrees, ruling out a
transitive $S_4$; (ii) the seam-lock gradient identity
$\partial\Ftil/\partial D_1 : \partial\Ftil/\partial S_1 = (3D_1^2+S_1):D_1$, the
first-stage normal ratio, distinguishes the seam-aligned pairing; (iii) the lifted
monodromy realises exactly the partition stabiliser. The group is dihedral, not
the Klein four group, because the product sits \emph{over} the seam
($P_2^2 = 3-P_1 C$): circling the seam swaps which product branch point is
encircled, which is the relation $\tau_+\leftrightarrow\tau_-$.

\section{The directive cover}
\label{sec:directive}

Solving $\Ftil$ for $D_1$ over $(S_1,S_2,P_2)$: now $P_2$ (hence $B$) is fixed, so
$A$ is no longer single-valued---it satisfies the \emph{$A$-block polynomial}
obtained by eliminating $D_1$ from the stage-two relation. Each admissible $A$
then yields $m_D$ values of $D_1$ through $A=3-D_1^{m_D}-D_1 S_1$.

\begin{result}[$A$-block count, $m_P=2$]
\label{res:ablock}
At $m_P=2$ the $A$-block polynomial has degree $m_D$ in $A$, verified symbolically
for $m_D\in\{2,3,4\}$. \Symbolic\ Hence the directive cover factors as $m_D$
blocks (the $A$-values) of $m_D$ sheets (the $D_1$-roots within each block), total
degree $m_D^2$, in agreement with Result~\ref{res:spectrum}.
\end{result}

\begin{result}[Directive-cover monodromy]
\label{res:directive}
At $m_P=2$ the directive cover is the full wreath $\Sym(m_D)\wr\Sym(m_D)$ of order
$(m_D!)^{m_D}\,m_D!$, imprimitive with the $A$-values as block system. \Measured\
at $m_D\in\{2,3\}$:
\[
  \text{cubic: } S_3\wr S_3 \ (\text{order } 1296),\qquad
  \text{keystone: } \Sym(2)\wr\Sym(2)\cong D_4 \ (\text{order } 8).
\]
\end{result}

\noindent For the cubic, the $9$ sheets split into $3$ $A$-blocks of $3$; nine
``within-block'' lassos realise $S_3$ in each block (the base $S_3^3$), and three
``$A$-branch'' lassos realise $S_3$ on the blocks. The generated group has order
$1296=6^3\cdot 6$, is transitive and imprimitive, with block kernel of order
$216=6^3$ and block action $S_3$---the full wreath. For the keystone, the $4$
sheets split into $2$ $A$-blocks of $2$; within-block lassos give $(0\,3),(1\,2)$
and a block-swap lasso (reached by looping $P_2$, not $S_1$) gives $(0\,1)(2\,3)$,
generating $D_4$. The cover is connected because the $A$-block discriminant
$\Delta = S_2^2(12-4P_2^2+S_2^2)$ is not a square in $\Q(S_2,P_2)$, so the two
$A$-values fuse.

\begin{remark}[The directive blow-up]
The contrast is the central asymmetry of the paper: when $m_D\neq m_P$, inverting
toward the directive is vastly larger than producing the product. For the cubic
the directive cover ($1296$) dwarfs the product cover ($8$).
\end{remark}

\section{The omega substrate}
\label{sec:omega}

The second-stage substrate $S_2$ enters $\Ftil$ only through the stage-two
coupling and never through the seam. We make this quantitative.

\begin{result}[Omega signature]
\label{res:omega}
$S_2$ satisfies $\partial A/\partial S_2 = 0$ and (writing $C$ for the stage-two
directive) $\partial C/\partial S_2 = 1$: it moves the product directive one-to-one
while leaving the seam invariant. No other coordinate does this---
$\partial A/\partial S_1 = -D_1\neq 0$ and $\partial A/\partial D_1 = -m_D
D_1^{m_D-1}-S_1\neq 0$. \Proven
\end{result}

\begin{result}[Branch authority]
\label{res:authority}
With $A$ held fixed, the product branches sit at $C = \pm 3/\sqrt{A}$, i.e.\ at
\[
  S_2^{(+)} = +\tfrac{3}{\sqrt A}-A \quad(\text{the }{+}P_1\text{ branch}),\qquad
  S_2^{(-)} = -\tfrac{3}{\sqrt A}-A \quad(\text{the }{-}P_1\text{ branch}).
\]
A loop of $S_2$ around $S_2^{(+)}$ swaps \emph{only} the $+P_1$ product pair
($(0\,1)$); a loop around $S_2^{(-)}$ swaps \emph{only} the $-P_1$ pair
($(2\,3)$). \Measured\ (residuals $\sim10^{-15}$).
\end{result}

\noindent Thus $S_2$ is an \emph{omega role}: a downstream branch-position
controller. It carries no seam structure and no directive structure; it sets where
the second-stage product branch lands, independently for each seam sign, and
steering it moves exactly the targeted product pair. Tying $S_2$ to another role
removes this: $S_2{=}0$ slaves the product branch to the seam, $S_2{=}S_1$ and
$S_2{=}D_1$ make the controlling knob co-move $A$ (contaminating the seam), and
$S_2{=}\text{const}$ freezes it. This also explains the substrate asymmetry that
runs through the construction: \emph{$S_1$ builds the seam, $S_2$ positions the
product}---first-stage substrate is upstream geometry, second-stage substrate is
the composition environment.

\section{Iterated composition and the recovery asymmetry}
\label{sec:iterate}

The glue iterates: stage three identifies $D_3:=P_2$, and so on.

\begin{result}[Recovery exponents]
\label{res:iterate}
For the three-stage composite of the cubic, the degree spectrum of
$(D_1,S_1,S_2,S_3,P_3)$ is $(27,9,6,4,8)$. \Symbolic\ In general the root-directive
recovery degree after $n$ stages is $m_D^{\,n}$ and the final-product production
degree is $m_P^{\,n}$:
\[
  \text{cubic backward: } 3,9,27,\dots=3^n;\qquad
  \text{cubic forward: } 2,4,8,\dots=2^n.
\]
\end{result}

\noindent Producing the product is monotonically cheaper than reconstructing the
directive, and the gap widens by a factor $m_D/m_P$ per stage. This is the
information-theoretic residue of the original ``category error'': inverting toward
the operator is not the same operation as producing the product run backwards---it
is a strictly higher-degree, more-branched cover, and the cost compounds through
every seam. No stage is ever cheaper than the previous: the seam never rejects.

\section{The abelian boundary on a phase interpretation}
\label{sec:abelian}

Covering-space monodromy and wave phase share their formalism: $\pi_1$ of a
punctured base mapping to automorphisms, and for cyclic (abelian) covers this is a
character into $U(1)$, with a lasso generating a phase $e^{2\pi i/k}$ and a branch
point as a phase singularity. It is therefore tempting to read the whole monodromy
as ``wave structure'' and the exponents $m_D^n, m_P^n$ as a frequency spectrum.
The following bounds that reading.

\begin{result}[Abelianisation]
\label{res:abelian}
For the directive cover of the cubic, $G=S_3\wr S_3$ has $|G|=1296$,
$|[G,G]|=324$, and abelianisation
\[
  G^{\mathrm{ab}} \;\cong\; (\Z/2)^2 \quad(\text{Klein four}).
\]
Hence every $U(1)$ character of $G$ takes values in $\{\pm 1\}$ (the parities of
the within-block and block permutations); an order-$3$ element maps to $+1$ under
every one of them. The cube-root phase $e^{2\pi i/3}$ is \emph{not} a character
value. \Measured
\end{result}

\noindent The cube roots are real but live elsewhere: they are the eigenvalues
$e^{\pm 2\pi i/3}$ of a $3$-cycle inside the $2$-dimensional irreducible
representation, which does not factor through $U(1)$. Thus the phase reading is
exact for the \emph{abelian shadow}---the $\pm1$ sign-phases, the
$m_P$-frequency---and categorically fails for the higher (cubic) structure, which
is the non-abelian part carrying the directive and the backward cost. In
particular, the forward/cheap structure is abelian and phase-like, while the
backward/expensive structure is non-abelian and phaseless; the boundary between
them is the boundary between $U(1)$ and what does not embed in it.

\section{Two-surface evidence: the keystone as falsification test}
\label{sec:keystone}

The keystone $m_D=m_P=2$ was chosen to be able to break the morphology: if the
directive cover threw up an $S_3$ or any $3$-cycle when both roles enter
quadratically, the law ``layer group $=\Sym(\text{entry degree})$'' would die. It
does not.

\begin{center}
\begin{tabular}{lcc}
\toprule
 & cubic $(m_D{=}3,m_P{=}2)$ & keystone $(m_D{=}m_P{=}2)$ \\
\midrule
degree spectrum $(D_1,S_1,S_2,P_2)$ & $(9,3,2,4)$ & $(4,2,2,4)$ \\
product cover & $D_4$ (order $8$) & $D_4$ (order $8$) \\
directive cover & $S_3\wr S_3$ (order $1296$) & $D_4$ (order $8$) \\
cube-root strata & yes (directive) & none \\
block structure & $3{\times}3$ / $2{\times}2$ & $2{\times}2$ on both \\
\bottomrule
\end{tabular}
\end{center}

\noindent Every entry matches the frozen prediction: both covers $D_4$, both
degree $4$, both imprimitive with $2$ blocks of $2$, and no $3$-cycle anywhere
(all branching is quadratic). The wreath-with-seam-imprimitivity structure now
stands on two surfaces, the second adversarial.

\section{Discussion and open problems}
\label{sec:discussion}

\paragraph{The seam as the fourth structural element.}\Interp\ The seam---the
directed identification $P_1=D_2$---is the load-bearing object: it is the
imprimitivity system of every cover, it carries the forward orientation of the
composition, and eliminating it \emph{is} enforcing that the product was produced
and passed on. Read as direction, it orients composition forward; iterated, it
orders the stages (an ordinal, not metric, arrow---there is no duration in the
construction); read as confirmation, it certifies the product. We emphasise this
is interpretation, and that the construction's confirmation is unconditional: the
seam has no failure branch (Result~\ref{res:iterate}), so ``feedback'' overreaches.

\paragraph{Beyond $m_P=2$.}\Conj\ The product cover is structurally
$\Sym(m_P)\wr\Sym(m_P)$ for any $m_P$ (two $m_P$-fold extractions), but this is
tested only at $m_P=2$. The decisive open test is $m_P=3$: does the product side
climb to $S_3\wr S_3$ as the directive side did? At $m_P\ge3$ the seam is a higher
root and the directive-cover $A$-block degree is no longer a clean power of $A$;
the directive-cover law $\Sym(m_D)\wr\Sym(m_D)$ is established (block structure) only
for $m_P=2$.

\paragraph{Fullness of the wreath.}\Conj\ The block structure of both covers is
symbolic for all $m_D$ (at $m_P=2$); that the monodromy realises the \emph{full}
wreath rather than a proper subgroup is computer-verified only at $m_D\in\{2,3\}$.
A proof that the lassos generate the full base and full block actions for all
$m_D$ would upgrade Results~\ref{res:product},\ref{res:directive} to theorems.

\paragraph{Universality across DBP surfaces.} All results are for the specific
family $D^{m_D}+DS+P^{m_P}-3$. Whether the seam-imprimitive wreath structure holds
for general separable DBP surfaces $h(P)=k(D,S)$ is open; separability of the
product is likely a genuine hypothesis (non-separable cross terms break the clean
radical tower).

\appendix
\section{Computational protocol}
\label{sec:methods}

Monodromy groups are obtained by numerical analytic continuation of the sheets
around based loops in the punctured base, under the following gates, all of which
must pass at every accepted continuation step:
\begin{enumerate}[leftmargin=1.4em,nosep]
\item \textbf{Root residual} on \emph{every} defining equation
(e.g.\ $|P_1^2-A|$ and $|P_2^2-(3-P_1C)|$), required $<10^{-9}$ (achieved
$\le4\times10^{-15}$).
\item \textbf{Bijective closure}: the nearest-neighbour sheet assignment must be a
bijection.
\item \textbf{Motion-ratio safety}: max sheet motion over min sheet separation is
tracked and kept small ($\sim10^{-5}$).
\item \textbf{Group from permutations}: the group is generated from the measured
permutations and its order computed by a permutation-group routine, never inferred.
\item \textbf{Byte-stability}: each run is repeated and the report's
content hash (timestamp excluded) compared; runs are byte-identical.
\end{enumerate}
Symbolic results (degree spectra, $A$-block degrees, the omega derivatives, the
abelianisation) are exact computer algebra. The product-cover artifact (loops,
generators, gates, generated group, partition-stabiliser check, and a test that
re-derives $D_4$ from the measured permutations) is archived separately.

\section*{Acknowledgements}
Computational consultation and verification were carried out in an interactive
session; all stated groups were re-derived from measured permutations under the
gates of Appendix~\ref{sec:methods}. \emph{[To be completed.]}

\end{document}
